\section{Toyota Production System (TPS) y Lean}

\begin{tcolorbox}[colback=red!5!white,colframe=red!75!black,title=Mapa Mental Rápido (Cópialo)]
\textbf{Filosofía Lean:} Eliminar desperdicios (Mudas) de forma implacable.\\
\textbf{Dos Pilares del TPS:} 
1. \textit{JIT (Just in Time):} Producir solo lo necesario, cuando es necesario (Kanban / Pull).
2. \textit{Jidoka:} Autonomatización con toque humano (Calidad en la fuente, parar la línea si hay error).
\end{tcolorbox}

\subsection{Los 7 Mudas (Desperdicios a evitar)}
Anota y memoriza esta lista, suelen cruzar conceptos en los V/F:
\begin{enumerate}
    \item \textbf{Sobreproducción:} El peor de todos (crea los demás). Hacer más de lo pedido.
    \item \textbf{Inventario:} Oculta los problemas bajo el agua (teoría del barco y las rocas).
    \item \textbf{Movimiento:} Personas moviéndose ineficientemente en su estación.
    \item \textbf{Transporte:} Mover materiales largas distancias innecesariamente.
    \item \textbf{Esperas:} Tiempos muertos de máquinas o personas.
    \item \textbf{Sobre-procesamiento:} Hacer cosas que el cliente no valora ni paga.
    \item \textbf{Defectos:} Reprocesos y fallas de calidad.
\end{enumerate}

\subsection{Conceptos y Herramientas (Vocabulario de Examen)}
\begin{itemize}[itemsep=1ex]
    \item \textbf{Genchi Genbutsu:} Ir a la fuente real del problema. No diagnosticar desde la oficina.
    \item \textbf{Kaizen:} Mejora continua con participación de todos, no solo de gerencia.
    \item \textbf{Heijunka:} Nivelación de la producción en volumen y mix (fabricar un poco de todo cada día, no por lotes masivos).
    \item \textbf{Kanban:} Tarjetas visuales que autorizan la producción. Son el nervio del sistema \textbf{PULL}.
    \item \textbf{Sistema Pull vs Push:} \textit{Pull} (jalar) reacciona a la demanda real. \textit{Push} (empujar) produce basado en pronósticos ciegos (como MRP).
    \item \textbf{Poka-Yoke:} Mecanismo a prueba de errores (ej. un conector que solo encaja de un lado). Apoya a Jidoka.
    \item \textbf{Andon:} Pizarra de luces para pedir ayuda o parar la línea.
    \item \textbf{5S:} Clasificar, ordenar, limpiar, estandarizar y sostener. Es disciplina visual y operacional del puesto de trabajo.
\end{itemize}

\begin{tcolorbox}[colback=blue!5!white,colframe=blue!70!black,title=Historia TPS que puede caer en V/F]
\textbf{Sakichi Toyota:} telar automático que se detenía al detectar falla. Es la raíz de \textbf{Jidoka}: calidad en la fuente y detención ante anomalías.\\
\textbf{Taiichi Ohno:} observa supermercados y desarrolla la lógica de reposición según consumo real: base del \textbf{Just-in-Time} y del sistema \textbf{Pull}.\\
\textbf{TPS no es solo técnica:} también es cultura de mejora continua, respeto, trabajo en equipo y resolución en la fuente.
\end{tcolorbox}

\subsection{Trampas Clásicas (V/F)}
\begin{itemize}
    \item \textit{``Lean significa cero inventario siempre.''} $\rightarrow$ \textbf{FALSO.} El inventario se minimiza, pero se mantiene un "buffer" estratégico para fluidez.
    \item \textit{``Jidoka es producir rápido y revisar la calidad al final.''} $\rightarrow$ \textbf{FALSO.} Es revisar en la fuente y parar inmediatamente.
    \item \textit{``En TPS se prefiere un lote grande para ahorrar setup.''} $\rightarrow$ \textbf{FALSO.} Se reduce el tiempo de setup (SMED) para poder usar lotes mínimos.
    \item \textit{``Heijunka significa producir siempre el mismo producto.''} $\rightarrow$ \textbf{FALSO.} Es nivelar volumen y mezcla para evitar variabilidad y sobrecargas.
    \item \textit{``Andon es solo un tablero informativo.''} $\rightarrow$ \textbf{FALSO.} Es un mecanismo para visibilizar anomalías y activar ayuda/detención.
\end{itemize}
